\src{187}
\result{COROLLAIRE}{} \uline{Pour qu'on ait} $\symup{H^{i}(X,F(-m))=0}$ \uline{pour} $\symup{m}$ \uline{grand, il
faut et il suffit que} $\symup{E^{n-i}(F)=0}$.

Rappelons que les $\symup{E^{j}(F)}$ ne dépendent pas en fait de l'immersion projective
envisagée. La condition du corollaire est de nature purement locale, par suite,
si elle est vérifiée pour $\symup{F}$, elle l'est pour tout faisceau localement
isomorphe à un faisceau $\symup{F^{n}}$. En particulier, si cette condition est vérifiée
pour $\symup{\Osh_X}$, elle l'est pour tout faisceau algébrique cohérent localement
libre. C'est par exemple le cas pour tout $\symup{i<n}$ si $\symup{X}$ est non singulière, pour
$\symup{i=0}$ si aucune composante de $\symup{X}$ n'est réduite à un seul point, pour $\symup{i=0,1}$
si $\symup{X}$ est normale et toutes ses composantes sont de dimension $\symup{>1}$ (cf.
[3]).\footnote{Le texte imprimé porte ici \fquote{$\symup{S}$ est normale} ; le paragraphe
ne fait intervenir que la variété $\symup{X}$.} Pour qu'il en soit ainsi pour tout
$\symup{i<k}$, il faut et il suffit par définition que les anneaux locaux
$\symup{\Osh_x}$ ($\symup{x\in X}$) soient de \fquote{codimension homologique} $\symup{\geq k}$ (cf [4]
pour cette notion). Si $\symup{k=n}$, cela signifie en vertu du théorème 3 bis que le
théorème de dualité est vrai pour $\symup{X}$, i.e. que (6.5) est un isomorphisme pour
tout $\symup{p}$ et tout $\symup{F}$. On peut donner de nombreuses conditions équivalentes sur
les anneaux locaux $\symup{\Osh_x}$ pour qu'il en soit ainsi (NAGATA), par exemple celle
de vérifier le théorème d'équidimensionnalité de Cohen-Macaulay. Il en est ainsi
par exemple si localement $\symup{X}$ est une \fquote{intersection complète} dans une
variété ambiante non singulière.

\fsection{\normalfont7}{\normalfont\uline{La dualité de Poincaré}.}

Soit $\symup{X}$ une variété projective non singulière de dimension $\symup{n}$. Alors
$\symup{H^{**}(X)=H^{*}(X,\underline{\Omega}_X^{*})}$ est une algèbre anticommutative bigraduée de dimension finie, que nous graduons
par le degré total, $\symup{H^{p,q}(X)=H^{p}(X,\underline{\Omega}_X^{q})}$ étant donc de degré
$\symup{p+q}$. Les degrés de $\symup{H^{*}(X)}$ sont compris entre $\symup{0}$ et $\symup{2n}$. En vertu du
théorème 2 et du corollaire au théorème 3, $\symup{H^{*}(X)}$ est une \fquote{algèbre de
Poincaré} de dimension $\symup{2n}$, i.e. $\symup{H^{2n}(X)}$ est muni d'un isomorphisme avec
le corps de base $\symup{\underline{k}}$, et le produit
$\symup{H^{m}(X)\times H^{2n-m}(X)\longrightarrow H^{2n}(X)=\underline{k}}$ est une dualité entre $\symup{H^{m}(X)}$ et $\symup{H^{2n-m}(X)}$. De plus, si $\symup{Y}$ est une
autre variété projective non singulière, la formule de Künneth pour les
faisceaux algébriques cohérents donne
\begin{equation}\tag{7.1}
  \symup{H^{*}(X\times Y)=H^{*}(X)\otimes H^{*}(Y)}
\end{equation}
isomorphisme compatible avec les structures d'algèbres de Poincaré. De plus,
$\symup{H^{*}(X)}$ est, en tant qu'algèbre commutative, un foncteur contravariant en
$\symup{X}$, un morphisme $\symup{f:Y\longrightarrow X}$ définissant de façon évidente un
homomorphisme d'algèbres graduées

